\begin{proof}[Proof of \cref{obj:lem:external-two-source-lower-certificate}]
Write \(\beta=1-2\alpha\), so \(0<\beta<1\), \(L_\alpha=1+\beta^2>1\), and both displayed certificate coefficients are positive once \(S_g>0\).

The interval \(\mathcal E\) is infinite, whereas the label set is finite. Thus some fiber of \(g\) contains three points \(e_1<e_2<e_3\). For any such triple, put
\[
D(e_1,e_2,e_3)
=\sqrt{(e_3-e_2)^2+(e_3-e_1)^2+(e_2-e_1)^2}.
\]
Overlap gives \(0<e_1/e_2<1\), and \(D(e_1,e_2,e_3)\ge e_3-e_2>0\). Consequently every ratio in the definition of \(S_g\) is positive and at most one. The set of ratios is nonempty, so
\[
0<S_g\le1,\qquad c_{\mathrm{tr}}(\alpha)>0,\qquad
c_{\mathrm{log}}(g,\alpha)>0.
\]
% lean: CausalSmith.PartialID.UnlinkedPropensityAte.fiberTripleSeparation_bounds

We first establish a uniform upper bound used to pass the score-log certificate through an infimum and a liminf. For finite measures on \(\mathcal V=[v_-,v_+]\), define
\[
d_{\mathcal V}(\mu,\nu)
=|\mu(\mathcal V)-\nu(\mathcal V)|
 +\int_{v_-}^{v_+}
   |\mu(\mathcal V\cap(-\infty,t])-\nu(\mathcal V\cap(-\infty,t])|\,dt.
\]
Let \(\lambda_{ar}(B)=P_{\mathrm{rel}}(A=a,R=r,Y\in B)\) and
\(\sigma_{ar}(B)=\int_{B\cap C_r}p_a(e)\,H(de)\). Using the empirical cell measures in \cref{obj:def:external-projection-handle}, set
\[
D_Y=\sum_{a,r}d_{[0,1]}(\widehat\lambda_{ar},\lambda_{ar}),
\qquad
D_E=\sum_{a,r}d_{\mathcal E}(\widehat\sigma_{ar},\sigma_{ar}),
\qquad
\Lambda_\varepsilon=\varepsilon^{-1}+\varepsilon^{-2}.
\]
For a retained sieve array, the triangle inequality places its outcome and score arrays within \(r_n+D_Y\) and \(r_m+D_E\), respectively, of the population arrays. The quantile-transport arm endpoints are Lipschitz in these distances with coefficient \(\Lambda_\varepsilon\): the reciprocal arm probabilities are bounded by \(\varepsilon^{-1}\), their variation across scores is bounded by \(\varepsilon^{-2}\), and integration of the corresponding distribution-function differences gives the two distances above. Hence each of the four arm endpoints differs from its population endpoint by at most
\[
\delta=\Lambda_\varepsilon(r_n+D_Y+r_m+D_E).
\]
The lower and upper ATE endpoints each combine two arm endpoints. Taking their convex hull and clipping to \([-1,1]\) therefore gives, on the projection branch of \cref{obj:def:external-projection-handle}, for its interval \(C_{n,m}\),
\[
\bigl(\operatorname{len}C_{n,m}
       -\operatorname{len}I(P_{\mathrm{rel}};H,g)\bigr)_+
\le 4\Lambda_\varepsilon(r_n+D_Y+r_m+D_E).
\]
The same inequality holds when the totalized interval in \cref{obj:def:external-projection-handle} is \(\{0\}\), since its length is zero. The empirical-CDF variance bound, integrated over each support interval and summed over the \(2J\) cells, yields
\[
\mathbb E D_Y\le\frac{2J}{\sqrt n},
\qquad
\mathbb E D_E\le\frac{J(2-2\varepsilon)}{\sqrt m}.
\]
The radii \(r_n=4J/(\alpha\sqrt n)\) and
\(r_m=2J(2-2\varepsilon)/(\alpha\sqrt m)\) make the projection interval honest: Markov’s inequality assigns probability at most \(\alpha/2\) to each failed empirical ball, and approximation by compatible rational arrays places the population sharp interval in the resulting closed projection hull on the complementary event. Thus \cref{obj:def:external-excess-risk} and the preceding bounds give
\[
\mathcal R^\star_{n,m,\alpha}(g)
\le 4\Lambda_\varepsilon
\left[
 \left(\frac{4J}{\alpha}+2J\right)n^{-1/2}
 +\left(\frac{2J(2-2\varepsilon)}{\alpha}
          +J(2-2\varepsilon)\right)m^{-1/2}
\right].
\]
In particular, with
\[
M=4\Lambda_\varepsilon
\left(\frac{4J}{\alpha}+2J+
      \frac{2J(2-2\varepsilon)}{\alpha}
      +J(2-2\varepsilon)\right)+1,
\]
we have \(T_{n,m,\alpha}(g)\le M\) for all \(n,m\ge1\), by the definition of \(b_{n,m}\) in \cref{obj:def:external-normalized-risk}.
% lean: CausalSmith.PartialID.UnlinkedPropensityAte.externalExcessRisk_normalized_upper

Fix an ordered triple \(e_1<e_2<e_3\) in one fiber, and define
\[
\mathbf h=(e_3-e_2,\;-(e_3-e_1),\;e_2-e_1),
\quad
u_m=\frac{\sqrt{\log L_\alpha}}
          {\sqrt3\,\|\mathbf h\|_2\sqrt m},
\quad
q=\frac{e_1+e_2+e_3}{3},
\quad
t=\frac12\left\{\frac{e_1}{3}
       +\min\!\left(\frac{e_1+e_2}{3},\frac{e_3}{3}\right)\right\},
\quad s=\frac tq.
\]
The identities \(\sum_i h_i=\sum_i e_i h_i=0\) show that
\[
H_0=\frac13\sum_{i=1}^3\delta_{e_i},
\qquad
H_m=\sum_{i=1}^3\left(\frac13+u_mh_i\right)\delta_{e_i}
\]
are probability laws with the same score mean \(q\) for all sufficiently large \(m\). The strict inequalities
\[
\frac{e_1}{3}<t<
\min\!\left(\frac{e_1+e_2}{3},\frac{e_3}{3}\right),
\qquad 0<s<1,
\]
hold by the choice of \(t\). Writing \(w_{i,m}=1/3+u_mh_i\) for \(i=1,2,3\), continuity and \(u_m\to0\) give, for all sufficiently large \(m\),
\[
w_{i,m}>0\quad(i=1,2,3),\qquad
e_1w_{1,m}<t<
\min\{e_1w_{1,m}+e_2w_{2,m},\,e_3w_{3,m}\}.
\] Under either score law, take \(Y_0=0\), an independent Bernoulli\((s)\) potential outcome \(Y_1\), score-based treatment assignment, and the fixed label of the triple. These constructions satisfy \cref{obj:def:external-law-class}; their released trial laws coincide, with treatment probability \(q\) and treated-success probability \(t\).

For score masses \(w_i\) in this neighborhood, the upper rearrangement of the treated successes fills score \(e_1\) and then score \(e_2\), while the lower rearrangement fits them at score \(e_3\). The endpoint formulas in \cref{obj:def:mean-endpoints} therefore reduce to
\[
U(w)=w_1+\frac{t-e_1w_1}{e_2},
\qquad L(w)=\frac{t}{e_3}.
\]
The control endpoints are zero. The two sharp intervals have common lower endpoint \(L(w)\), and their upper endpoints differ by
\[
\Delta_m
=U\!\left(\tfrac13+u_mh_1,\tfrac13+u_mh_2,\tfrac13+u_mh_3\right)
 -U\!\left(\tfrac13,\tfrac13,\tfrac13\right)
=u_m(e_3-e_2)\left(1-\frac{e_1}{e_2}\right)>0.
\]
By the endpoint-attainment assertion of \cref{obj:lem:fixed-released-law-baseline}, choose causal populations attaining the lower endpoint under \(H_0\) and the upper endpoint under \(H_m\), while preserving their respective released laws. If \(W_0\) denotes the sharp interval length under \(H_0\), their ATEs differ by \(W_0+\Delta_m\).

Only the score logs distinguish these endpoint-attaining experiments. Direct calculation gives
\[
\chi^2(H_m,H_0)
=3u_m^2\|\mathbf h\|_2^2
=\frac{\log L_\alpha}{m}.
\]
For product laws, \(1+\chi^2\) raises to the sample-size power; Cauchy--Schwarz then bounds total variation by half the square root of \(\chi^2\). Using \((1+x/m)^m\le e^x\), the joint experiments consequently satisfy, uniformly in \(n\),
\[
\operatorname{TV}(\mathbb Q_0^{n,m},\mathbb Q_m^{n,m})
\le\frac12
 \sqrt{\left(1+\frac{\log L_\alpha}{m}\right)^m-1}
\le\frac12\sqrt{L_\alpha-1}
=\frac{\beta}{2}.
\]
For any honest interval, transfer coverage of the second ATE to the first experiment and intersect it with coverage of the first ATE. The intersection has first-experiment probability at least
\(1-2\alpha-\operatorname{TV}\ge\beta/2\). On it the interval length is at least \(W_0+\Delta_m\). Taking the supremum over laws and the infimum over honest procedures gives, for every \(n\ge1\) and all sufficiently large \(m\),
\[
\mathcal R^\star_{n,m,\alpha}(g)
\ge\frac{\beta\Delta_m}{2}
=\frac{\beta\sqrt{\log L_\alpha}}{2\sqrt3\,\sqrt m}
  \frac{(e_3-e_2)(1-e_1/e_2)}
       {D(e_1,e_2,e_3)}.
\]
% lean: CausalSmith.PartialID.UnlinkedPropensityAte.scoreLogFixedTriple_eventually_externalExcessRisk_lower

To pass from a fixed triple to \(S_g\), take any number below
\(\beta\sqrt{\log L_\alpha}S_g/(2\sqrt3)\). The definition of the supremum supplies a triple whose displayed coefficient exceeds that number. Its eventual lower bound holds uniformly over \(n\ge1\). The infimum over \(n\) is finite: it is bounded below by this eventual certificate and above by the choice \(n=m\), for which \(b_{m,m}=\sqrt m/2\) and the bound \(T_{m,m,\alpha}(g)\le M\) gives
\(\sqrt m\,\mathcal R^\star_{m,m,\alpha}(g)\le2M\). Taking the liminf and then letting the chosen number increase to the supremum proves
\[
\liminf_{m\to\infty}\inf_{n\ge1}
 \sqrt m\,\mathcal R^\star_{n,m,\alpha}(g)
\ge c_{\mathrm{log}}(g,\alpha).
\]
% lean: CausalSmith.PartialID.UnlinkedPropensityAte.scoreLog_logCertificate_liminf

For the trial certificate, take \(H=\delta_{1/2}\), which is supported on \(\mathcal E\) by \cref{obj:ass:overlap}. Let \(P_0\) have \(Y_0=0\) and an independent Bernoulli\((1/2)\) outcome \(Y_1\), and let \(P_1\) replace its success probability by \(1/2+v_n\), where
\[
v_n=\frac12\sqrt{\frac{\log L_\alpha}{n}}.
\]
Both laws use treatment probability \(E=1/2\), consistency, and the fixed release \(g(1/2)\). Since \(0<\beta<1\) and \(n\ge1\), \(0<v_n<1/2\); both pairs belong to \cref{obj:def:external-law-class}. The fixed score makes the reciprocal treatment weights constant, so \cref{obj:def:mean-endpoints,obj:def:sharp-ate-set} give singleton sharp intervals at \(1/2\) and \(1/2+v_n\).

For a latent Bernoulli outcome, the chi-squared divergence between these two success probabilities is \(4v_n^2\). Adding the common randomized arm, releasing the trial row, and adding the common score log cannot increase total variation. Hence
\[
\operatorname{TV}(\mathbb Q_0^{n,m},\mathbb Q_1^{n,m})
\le\frac12\sqrt{(1+4v_n^2)^n-1}
\le\frac12\sqrt{e^{\log L_\alpha}-1}
=\frac{\beta}{2},
\]
where the middle inequality uses
\((1+\log L_\alpha/n)^n\le L_\alpha\). Under the first experiment, honesty and transfer of the second coverage event put both ATEs in any honest interval with probability at least \(\beta/2\). Its length on that event is at least \(v_n\), while the first sharp interval has length zero. Thus every honest procedure has worst-case expected excess length at least
\[
\frac{\beta v_n}{2}
=\frac{\beta\sqrt{\log L_\alpha}}{4\sqrt n}
=\frac{c_{\mathrm{tr}}(\alpha)}{\sqrt n}.
\]
The full interval \([-1,1]\) is honest, so the infimum defining the minimax risk is over a nonempty set; every candidate satisfies the displayed bound. This proves the finite-sample assertion.
% lean: CausalSmith.PartialID.UnlinkedPropensityAte.trialCenter_honest_expectedLength

Finally, for positive \(n,m\), multiplying the trial bound by \(b_{n,m}\) gives
\[
T_{n,m,\alpha}(g)
\ge c_{\mathrm{tr}}(\alpha)
 \frac{\sqrt m}{\sqrt n+\sqrt m}.
\]
If \(n,m\to\infty\) with \(n/m\to\rho>0\), the fraction tends to
\(1/(1+\sqrt\rho)\). For the log bound, let any
\(\gamma<c_{\mathrm{log}}(g,\alpha)\) be fixed. The preceding liminf inequality implies, eventually in \(m\), uniformly over \(n\ge1\),
\(\mathcal R^\star_{n,m,\alpha}(g)\ge\gamma/\sqrt m\). Hence eventually
\[
T_{n,m,\alpha}(g)
\ge \gamma\frac{\sqrt n}{\sqrt n+\sqrt m}
\longrightarrow
\gamma\frac{\sqrt\rho}{1+\sqrt\rho}.
\]
The uniform bound \(T_{n,m,\alpha}(g)\le M\) makes the normalized liminf finite. Letting \(\gamma\) increase to \(c_{\mathrm{log}}(g,\alpha)\) and combining the two lower bounds yields
\[
\liminf_{\substack{n,m\to\infty\\ n/m\to\rho}}
T_{n,m,\alpha}(g)
\ge
\max\left\{
\frac{c_{\mathrm{tr}}(\alpha)}{1+\sqrt\rho},
\frac{c_{\mathrm{log}}(g,\alpha)\sqrt\rho}{1+\sqrt\rho}
\right\}>0,
\]
as required.
% lean: CausalSmith.PartialID.UnlinkedPropensityAte.externalExcessRisk_lower_certificates
\end{proof}
